% Page 059 — page imprimée 56
% Georges ROTH — suite

Et toute la famille se retrouva à Meyrueis dans une maison qui donnait sur l’Eglise, à côté de
celle d’Huguette Poujol. Elle appartenait au pasteur \textbf{Etienne Giran}*, collègue de Georges
Roth qui avait les mêmes convictions libérales. Rencontré dans un train et apprenant qu’il
était nommé à Meyrueis il lui avait dit :
« J’ai là-bas une maison, occupez-la tant que vous voudrez ». Et il lui fit un mot pour qu’on
lui remette les clés.

Il devait bien y avoir parmi ces 1600 jeunes mobilisés dans les Chantiers de Jeunesse de
Meyrueis et des environs quelques protestants confrontés à la propagande de Vichy : retour à
la terre, recherche des vraies valeurs de la France profonde, vénération du Maréchal, du
drapeau ...

Georges Roth ne tint pas le langage convenu.
Il retrouva quelques alsaciens-lorrains qui avaient été démobilisés, ou se trouvaient à
l’Intérieur et ne rentrèrent pas en Alsace Moselle redevenue Allemande. Plusieurs jeunes,
dont ces alsaciens, demandèrent à l’aumônier protestant d’aménager un local dont ils
pourraient disposer. Un décor d’Alsace fut reconstitué avec alsacienne et alsacien en costumes
traditionnels, les grandes dates de l’histoire de l’Alsace mais surtout un verset biblique
explicite : « \textbf{En ce jour-là, dit l’Eternel des Armées, je briserai ton joug sur ton cou}\\
\textbf{Je romprai tes liens et des étrangers ne t’assujettiront plus} ».
\hfill Jérémie 30/8

Le pasteur Roth tenait en public et même dans des repas officiels, des propos subversifs
tendant à mettre en doute l’obéissance aux ordres venus de Vichy, prônant la Révolution
Nationale, propos qui furent rapportés aux chefs de groupement ou de région.

Il fut dénoncé par ses supérieurs, comme personnage néfaste, sapant l’autorité des chefs,
dénonciation envoyée au Pasteur Boegner, Président de la Fédération Protestante de France,
mais aussi probablement au Général de la Porte du Theil, commandant les Chantiers de
Jeunesse ; il semble que l’affaire remonta jusqu’à Pétain qui l’aurait stoppée. Peut-être...

Toujours est-il que la police allemande débarqua un beau matin à Meyrueis pour interroger et
peut-être arrêter G. Roth. Dans un premier interrogatoire l’officier apprit que Roth avait fait
ses études de théologie protestante en partie à Montpellier puis à Strasbourg comme lui-même,
logeant peut être dans les mêmes bâtiments du quai St. Thomas qu’on appelle toujours le Stift.
Il aurait ouvert la fenêtre en lui disant de s’enfuir.

Georges Roth ne se le fit pas dire deux fois et gagna le maquis. Mais pas celui de Bir Hakeim
qui sera anéanti à la Parade : il avait déjà eu des contacts avec les chefs mais avait renoncé à
les rejoindre, vu leur amateurisme et leur manque de rigueur et de prudence. Il leur en avait
fait la remarque en tant qu’officier mais n’avait pas été écouté.
Il rejoignit un des maquis de l’Aigoual, probablement à l’Espérou, qui disposait d’une
meilleure infrastructure avec des groupes disciplinés et mieux entraînés.

Après la Libération Georges Roth retrouva sa paroisse du Havre, ville à demi - détruite, puis
fut aumônier « cette fois pour de bon ! » dans l’armée d’occupation en Autriche avant d’être
appelé à l’Eglise Wallonne d’Amsterdam (comme son ami Etienne Giran qui, déporté à
Buchenwald , n’en revint jamais).

\bigskip

\noindent * Etienne Giran : [la note de bas de page n’apparaît pas sur cette page.]
